% ==============================================================================
% РАЗДЕЛ 2.3: Анализ применения преобразования Фурье
% Партия B03-020 (DISS-005-B0862 .. DISS-005-B0891, 30 блоков)
% ==============================================================================

% DISS-005-B0862
\subsection{Анализ применения преобразования Фурье}
\label{sec:ch02_fourier_theory}

% DISS-005-B0863
Современные системы релейной защиты и автоматики (РЗА) работают на микропроцессорной основе, что предполагает дискретизацию аналоговых сигналов. Для анализа спектральных характеристик токов и напряжений, необходимых для принятия решений в БАВР, применяется дискретное преобразование Фурье (ДПФ). Однако ДПФ, работая с ограниченным по времени набором данных, подвержен эффекту спектрального растекания. Для минимизации этого эффекта применяются оконные функции. Они позволяют улучшить точность оценки амплитуд и фаз гармоник. Поэтому рассмотрение различных оконных функций и их влияния на точность анализа переходных процессов в БАВР является очень важным [36, 37].

% DISS-005-B0864
Оценка оконных функций проводилась в MATLAB путем перебора различных вариантов ширины окна ($L$) и постоянной затухания ($Ta$). Для некоторых оконных функций, как например для треугольной актуально лишь использование окна минимальной ширины в 2 периода промышленной частоты. При частоте дискретизации 96 выборок на период, угол короткого замыкания задавался равным 90 градусам для обеспечения максимального броска тока.

% DISS-005-B0865
Использовались следующие параметры:

% DISS-005-B0866
\paragraph{Время вхождения в зону}

% DISS-005-B0867
$t_{\text{зона}}$~-- это временной интервал, отсчитываемый с момента возникновения КЗ до момента, начиная с которого отклонение расчетного действующего значения тока от истинного действующего значения не превышает заданной величины ($\pm 1\%$). Используется для оценки времени наступления установившегося значения и прекращению колебательных процессов.

% DISS-005-B0868
\paragraph{Время достижения уставки}

% DISS-005-B0869
$t_1$~-- время первого достижения истинного действующего значения. Используется для оценки быстродействия различных органов, предназначенных для защиты от короткого замыкания.

% DISS-005-B0870
\paragraph{Время достижения первого максимума}

% DISS-005-B0871
$t_{1\max}$~-- время достижения первого максимума. Используется для оценки времени возникновения первого броска.

% DISS-005-B0872
\paragraph{Максимальное отклонение амплитуды}

% DISS-005-B0873
$\sigma_{\max}$~-- это максимальное отклонение расчетного действующего значения тока в большую сторону от истинного действующего значения, выраженное в процентах. Используется для оценки возможного перерегулирования в следствии чего может возникать ложное срабатывание токовых защит.

% DISS-005-B0874
% DISS-005-B0875
% DISS-005-B0876
\begin{equation}
\label{eq:ch02_max_amplitude_deviation}
\sigma_{\text{макс}} = 100 \cdot \left( \frac{I_{\text{макс}} - I_{\text{уст}}}{I_{\text{уст}}} \right)
\end{equation}

% DISS-005-B0877
Данные три параметра в рамках стандартной логики функционирования БАВР актуально для органа реле направления мощности (РНМ) и максимальной токовой защиты (МТЗ~-- 50/51).

% DISS-005-B0878
По результатам проведённых измерений (часть из которых представлена в таблица 1), а также ряда дополнительных, можно сделать следующие выводы

% DISS-005-B0879
\paragraph{Оптимальные окна для быстрого отслеживания:}
\begin{itemize}[noitemsep,topsep=2pt]
% DISS-005-B0880
\item Минимальное время достижения истинного значения и первого максимума тока обеспечивают прямоугольное окно и окно Кайзера (1 период, $Ta=0{,}1$~с) при использовании сигнала от ТТ ($t_1 \approx 13$~мс, $t_{1\max} \approx 25$~мс).
% DISS-005-B0881
\item Для сигнала производной тока от КР прямоугольное окно (1 период, $Ta=0{,}1$~с) сокращает зону отклонения до 18,5~мс (в 10 раз быстрее ТТ) с погрешностью $\le 1\%$.
\end{itemize}

% DISS-005-B0882
\paragraph{Управление переходными процессами:}
% DISS-005-B0883
Для апериодических составляющих предпочтительны:
\begin{itemize}[noitemsep,topsep=2pt]
% DISS-005-B0884
\item Треугольное окно (ТТ, $Ta=0{,}05 \div 0{,}1$~с)~--- баланс скорости и стабильности.
% DISS-005-B0885
\item Прямоугольное окно (КР, $Ta=0{,}05$~с)~--- минимальные $t_{1\max}=30$~мс, $t_{\text{зона}}=19$~мс.
% DISS-005-B0886
\item При резких скачках эффективно треугольное окно, при плавных изменениях~--- окно Кайзера.
\end{itemize}

% DISS-005-B0887
\paragraph{Подавление гармоник:}
\begin{itemize}[noitemsep,topsep=2pt]
% DISS-005-B0888
\item Прямоугольное окно минимизирует вторую гармонику (ТТ: 28,39\%, КР: 13,88\%).
% DISS-005-B0889
\item Узкие окна снижают гармоники, но усиливают спектральное растекание.
\end{itemize}

% DISS-005-B0890
\paragraph{Итог:} Выбор окна зависит от приоритета (скорость/точность/подавление помех). Прямоугольное окно с шириной в 1 период промышленной частоты оптимально для быстрого анализа и подавления гармоник с повышением устойчивость, что как раз особо важно для БАВР. Треугольное и Кайзера~-- для компромисса в динамических режимах, хоть и редко используются на практике.

% DISS-005-B0891
\paragraph{Таблица 1. Сравнение оконных функций}
